\section{Weighted-Wiener stability of the isolated resonance mode}
\label{sec:wiener-resonance-stability}
The finite-Fourier theorem above has an explicit infinite-support stability
region.  This is useful because finite support and arbitrary smooth support
are qualitatively different: smooth exact coboundary cancellations exist, but
they cannot approach a fixed finite polynomial too closely in a suitable
weighted Wiener topology while preserving zero variance.

For a Fourier series $f$ put
\[
 M_0(f)=\sum_{v\in\mathbb Z^2}\|v\|_1|\widehat f(v)|,
 \qquad
 M_A(f)=\sum_{v\in\mathbb Z^2}\|Av\|_1|\widehat f(v)|.
 \tag{W.1}\label{eq:wiener-moments}
\]
Constants are invisible to these seminorms and to the Poisson bracket.  Write
\[
 B_A(f,g)=\frac{\{f,\mathcal C_A g\}}{4\pi},
 \qquad h_g=B_A(g,g).
\]
The Fourier formula is
\[
 \widehat{B_A(f,g)}(k)
 =-\pi\!\sum_{u+Av=k}\det(u,Av)\widehat f(u)\widehat g(v).
 \tag{W.2}\label{eq:wiener-bracket}
\]
Since $|\det(u,w)|\leq\|u\|_1\|w\|_1$, absolute summation gives
\begin{equation}\label{eq:wiener-bilinear}
 \|\widehat{B_A(f,g)}\|_{\ell^1}
 \leq \pi M_0(f)M_A(g).
\end{equation}
Thus all orbit sums are absolutely convergent whenever the two weighted
moments are finite.

\begin{theorem}[Explicit weighted-Wiener stability]
\label{thm:wiener-resonance-stability}
Let $p$ be a nonconstant real trigonometric polynomial.  Choose the isolated
finite-support witness $u,v$ supplied by Lemma~\ref{lem:temporal-support},
for lag one, and put
\[
 d=|\det(u,Av)|,\qquad
 D=d\,|\widehat p(u)\widehat p(v)|>0.
\]
For a real weighted-Wiener perturbation $g=p+r$, define
\[
 P_0=M_0(p),\quad P_A=M_A(p),\quad
 \varepsilon_0=M_0(r),\quad \varepsilon_A=M_A(r)
\]
and
\[
 E=P_0\varepsilon_A+\varepsilon_0P_A
       +\varepsilon_0\varepsilon_A.
\]
If $E<D$, then
\begin{equation}\label{eq:wiener-gamma-lower}
 \Gamma_A(g)\geq2\pi^2(D-E)^2>0.
\end{equation}
\end{theorem}

\begin{proof}
Bilinearity and \eqref{eq:wiener-bilinear} give
\[
 \|\widehat{h_g-h_p}\|_{\ell^1}
 \leq\pi E.
\]
If $k=u+Av$, the isolated-mode theorem says that the oriented $A$-orbit of
$k$ contains exactly one Fourier coefficient of $h_p$, with modulus $\pi D$.
Hence the same orbit sum for $h_g$ has modulus at least $\pi(D-E)$.
For real $g$ the distinct orbit of $-k$ has the conjugate sum, so those two
orbits alone give \eqref{eq:wiener-gamma-lower}.
\end{proof}

A one-parameter radius follows immediately.  If
\[
 P=P_0+P_A,
 \qquad \varepsilon=\max\{M_0(r),M_A(r)\},
\]
then $E\leq P\varepsilon+\varepsilon^2$.  Therefore every perturbation with
\begin{equation}\label{eq:wiener-radius}
 \varepsilon<\varepsilon_*(p;u,v)
 :=\frac{\sqrt{P^2+4D}-P}{2}
\end{equation}
has positive resonance variance.  In particular every nonconstant real
trigonometric polynomial is an interior point of the nondegenerate set in
this weighted Wiener topology.

For analytic tails the criterion is fully explicit.  If
$|\widehat r(v)|\leq C e^{-\rho\|v\|_1}$ and the tail is supported on
$\|v\|_1>R$, then
\[
 M_0(r)\leq4C\sum_{n>R}n^2e^{-\rho n},
 \qquad M_A(r)\leq\|A\|_{1\to1}M_0(r).
 \tag{W.3}\label{eq:wiener-analytic-tail}
\]
Substitution into the theorem gives a computable analytic-tail condition.
This proves nondegeneracy on an infinite-dimensional open neighbourhood of
every finite polynomial; it does not prove that every nonconstant analytic
observable lies in one of these neighbourhoods.  The latter global coverage
question remains open.
